% =====================================================================
\section{The NMDA receptor}
The NMDA receptor opens only when two conditions coincide: the transmitter has arrived \emph{and} the cell is already excited. Calcium enters the cell through it. It therefore serves not only as an input but as an activity sensor: from its operation the cell ``learns'' that the network needs it, and maintains production of its proteins, both working and protective.
\facts
\begin{itemize}
\item Ketamine blocks the NMDA receptor. In healthy volunteers it produces states resembling the symptoms of schizophrenia and impairs attention and memory.\tP
\item A paradox: an excitatory receptor is blocked, yet cortical activity rises. The accepted explanation is that receptors on inhibitory cells suffer more, inhibition weakens, and pyramidal cells are \emph{disinhibited}.\tU
\item In cell culture ketamine halves the amount of GAD67 in about 4 hours and of parvalbumin in about 8. The effect is reversible and the cells do not die.\tS
\item In mice a single dose of ketamine leaves no trace, whereas two doses on two consecutive days cause a persistent loss of parvalbumin and GAD67, with slow and incomplete recovery by day 10.\tS
\item Blocking the receptor in newborn rodents halves, within a day, the activity of the enzyme that makes the cell's principal antioxidant.\tS\ Synaptic activity switches on genes for protective enzymes in neurons, among them the one that repairs damaged defence (section~9).\tP
\item Oxidants reversibly suppress the NMDA receptor itself.\tS
\end{itemize}
\begin{ineq}
Three facts compose a closed loop: the receptor maintains the defence; without defence oxidants rise; oxidants suppress the receptor. The two-dose experiment shows a threshold not only in the strength of the insult but in its \emph{duration}: the system has a memory at least a day long.
\end{ineq}
